\chapter*{前言：这本书不教你写代码}
\addcontentsline{toc}{chapter}{前言：这本书不教你写代码}
\markboth{}{前言}

过去，使用计算方法解决一个真实问题，往往意味着从学习编程开始。我们需要熟悉一门语言，理解各种软件库，逐行编写代码，再花费大量时间处理环境配置、语法错误和接口变化。对于许多学生、研究人员和工程师而言，真正想解决的问题还没有开始，时间已经消耗在实现细节之中。

生成式 AI 正在改变这种工作方式。今天，只要能够比较清楚地描述任务，AI 就可以帮助读取数据、生成代码、解释报错、构建模型、绘制图表，甚至整理一份初步报告。许多过去需要数天完成的技术实现，现在可能在数小时内得到一个看起来相当完整的结果。

但这并不意味着问题已经被解决。

AI 可以迅速给出答案，却不知道这个答案是否值得相信；它可以生成一套完整的分析流程，却未必理解数据是如何产生的；它可以选择一个复杂模型并报告漂亮的指标，却可能忽略数据泄露、错误标签、领域约束或现实中无法接受的错误成本。AI 提高了我们得到答案的速度，也放大了我们判断答案的责任。

本书有一个贯穿始终的核心命题：

\begin{center}
\textbf{AI 负责生成候选答案，人负责定义正确答案的标准。}
\end{center}

这是本书理解 AI 协作的基本出发点。AI 的优势在于快速生成方案、代码、解释和可能路径；人的责任则在于提出问题、约束边界、设定评价标准，并判断这些候选答案是否真正符合现实需求。

围绕这一核心命题，本书展开一条真实问题求解的协作主线：

\begin{center}
\textbf{如何提问（Prompting）}
$\longrightarrow$
\textbf{如何建模（Modeling）}
$\longrightarrow$
\textbf{如何判断（Evaluation）}
\end{center}

首先，我们需要把现实需求转化为 AI 能够理解和执行的提示词；随后，与 AI 一起把提示词落实为数据、模型和计算流程；最后，对生成的答案进行验证，判断它是否正确、可信并适用于现实。这三个阶段不是相互独立的技术步骤，而是一条反复迭代的协作链条：判断中发现的问题，会推动我们重新提问和重新建模。

\section*{一个尚待补充的开篇案例}

\begin{quote}
\textbf{【案例占位】}

这里将放入一个完整的交通运输工程案例：一名学习者面对一份真实数据和一个尚未被清楚定义的问题，借助 AI 快速完成数据处理、建模、可视化和报告，并得到一个表现优异但实际错误的结果。案例需要展示错误是如何产生的、为什么不容易被发现，以及领域知识和验证过程如何揭示问题。

\textbf{【待确定内容】} 数据来源、研究任务、AI 协作过程、表面上优秀的结果、隐藏错误及其现实后果。
\end{quote}

这个案例最终要说明的并不是“AI 会犯错”这样一个简单结论，而是一种更重要的工作原则：AI 能让我们更快到达一个答案，但不能保证那个答案是对的。过去，技术实现的困难会迫使我们缓慢前进；现在，生成答案变得容易，真正困难的部分则转移到了问题定义、约束表达和结果验证。

\section*{这本书为什么不教你写代码}

本书书名中的“提示即代码”，并不是说自然语言已经完全取代了程序，也不是说理解代码从此不再重要。它强调的是一种新的协作界面：我们越来越多地通过描述目标、提供上下文、声明约束和规定验收标准，来组织 AI 完成技术工作。

在这种方式下，人不必从零开始编写每一行代码，但必须知道自己正在要求 AI 做什么。一个有效的 Prompt 不是一句随意的问题，而是一份微型任务说明书。它需要回答：问题是什么，数据代表什么，哪些条件不能被破坏，结果应该如何检查，以及什么样的答案才算合格。

因此，本书不会按照传统编程教材的方式，从变量、循环和函数开始逐项讲解。我们更关心代码之前和代码之后的事情：

\begin{itemize}
  \item \textbf{如何提问：}在代码生成之前，把模糊的现实需求转化为明确、可计算的任务；
  \item \textbf{如何建模：}在 AI 执行过程中，把数据、专业背景、业务规则与必要约束落实到模型中；
  \item \textbf{如何判断：}在代码运行之后，检验输出是否正确、可信并适用于真实场景。
\end{itemize}

代码仍然重要，但它不再是全部。能够生成代码，也不等于能够解决问题。

\section*{Vibe Coding到底是什么}

Vibe Coding 常被理解为：使用者通过自然语言描述意图，让 AI 持续生成和修改代码，而不再亲自完成每一个实现细节。这种方式显著降低了从想法到原型的成本，也让更多人能够快速尝试过去难以完成的技术任务。

然而，Vibe Coding 不应被理解为“不需要懂代码”，更不意味着“不需要理解问题”。AI 无法自动判断一个研究问题是否合理，不知道交通数据中的异常值代表设备故障、拥堵还是事故，也不会天然发现所有逻辑漏洞、数据泄露和评价偏差。它擅长执行被表达出来的任务，却无法替人补齐所有未被表达的专业知识。

人的角色因此正在发生变化：

\begin{center}
代码的逐行实现
$\longrightarrow$
目标定义、约束控制与结果判断
\end{center}

这并不是人的价值被削弱，而是人的价值从“如何实现”更多地转向“为什么这样做”“什么不能做”和“如何证明做对了”。当 AI 可以快速生成多个候选方案时，提出正确问题、识别关键约束和判断证据质量，反而成为更稀缺的能力。

\section*{Prompt as Code意味着什么}

传统代码要求计算机严格执行一组明确指令。Prompt 则通过自然语言向 AI 表达意图，由 AI 生成具体实现。两者并不相同，但一个高质量 Prompt 与高质量代码具有相似特征：目标明确、边界清楚、结构合理、结果可验证。

把 Prompt 当作 Code，意味着我们不能把提示仅仅看作一句“向 AI 提问的话”。它应当像代码一样被设计、检查和迭代：

\begin{enumerate}
  \item 明确任务目标，避免让 AI 猜测真正需求；
  \item 提供必要上下文，让 AI 理解数据与业务场景；
  \item 声明专业约束，防止技术上可行但现实中错误的操作；
  \item 规定输出形式和验收标准，使结果可以被系统检查；
  \item 保留协作过程与关键决策，使分析能够被复核和改进。
\end{enumerate}

提示不是魔法咒语。真正有效的提示来自对问题本身的理解。

\section*{本书的逻辑主线}

当 AI 已经能够写代码之后，人的核心价值不再只是完成实现，而是组织一次可信的问题求解过程。本书因此不以工具或算法为中心，而以“如何提问、如何建模、如何判断”为逻辑主线。

\begin{description}
  \item[如何提问（Prompting）] 提问决定了 AI 正在解决什么问题。我们需要明确目标、对象、范围、约束与验收标准，把含混的现实需求整理成可以执行和验证的任务。好的 Prompt 不是追求某种固定句式，而是准确表达对问题的理解。
  \item[如何建模（Modeling）] 建模决定了现实问题如何被数据和计算表示。我们需要理解数据的来源与含义，完成必要的数据处理，选择合适的方法，并将领域知识、业务规则和风险边界传递给 AI。模型越复杂，越需要清楚它建立在什么假设之上。
  \item[如何判断（Evaluation）] 判断决定了 AI 给出的答案能否被采用。我们需要使用基准比较、数量级检查、极端情景、可视化和证据链验证输出，同时识别数据泄露、评价偏差与不合理结论。一个漂亮的指标，只是判断的起点。
\end{description}

全书各部分围绕这三个问题逐步推进：第 1 至第 4 章帮助读者理解 AI、学习有效沟通，并把现实需求转化为清楚的问题；第 5 至第 11 章讨论如何准备数据、构造特征，并与 AI 协作完成不同类型的建模任务；第 12 至第 14 章重点讨论如何验证结果、识别错误与能力边界；第 15 至第 18 章则综合运用 Prompting、Modeling 与 Evaluation，将经过检验的方案推进到完整项目和真实应用中。

这条主线并非单向流程。一次评价可能暴露模型假设的问题，一次建模尝试也可能说明最初的问题定义并不准确。可靠的 AI 协作，正是在 Prompting、Modeling 与 Evaluation 之间不断往返，直到问题、方法与证据能够相互支持。

本书主要使用交通运输工程领域的问题作为案例来源。这些问题通常包含复杂的时空结构、明确的现实约束和不可忽视的错误成本，非常适合展示 AI 协作的价值与边界。但书中讨论的方法并不限于交通领域，也可以迁移到其他需要数据分析、建模和专业判断的真实问题中。

\section*{如何使用本书}

本书不是一本适合从头到尾被动阅读的知识手册。更合适的使用方式，是带着一个真实问题阅读，并在每个阶段与 AI 实际协作。

\begin{description}
  \item[本科生] 可以把本书作为“AI 辅助完成课程项目”的指南，重点学习如何描述问题、检查代码和解释结果。
  \item[研究生] 可以把本书作为“AI 辅助科研建模与实验设计”的手册，重点关注数据泄露、方法选择、验证流程和证据质量。
  \item[工程师] 可以把本书作为快速完成业务分析原型的工具书，重点关注业务约束、部署边界、风险控制和成果表达。
\end{description}

阅读过程中，不必追求记住每一种工具或每一个 Prompt 模板。AI 工具会快速变化，具体界面和能力也会不断更新。更值得长期掌握的是一套稳定的方法：先定义问题，再组织 AI 执行；先明确判断标准，再接受 AI 给出的答案。

\section*{最后，责任仍然属于人}

AI 可以参与分析，但不能承担分析的责任。无论代码、图表或文字由谁生成，最终采用这些结果并据此作出判断的人，都必须对其正确性和后果负责。

这也是本书最希望传递的态度：不要因为 AI 能够快速回答，就轻易放弃思考；也不要因为 AI 会犯错，就拒绝使用它。真正有价值的能力，是知道什么时候应该相信，什么时候应该怀疑，什么时候应该降低 AI 的角色，以及什么时候必须由人作出最终决定。

当代码生成不再是最大的障碍，我们终于可以把更多注意力放回真正的问题上。

\vspace{2em}
\begin{flushright}
编者\\
\today
\end{flushright}
